\section{Conclusion}
\label{sec:conclusion}

With a few target labels per emotion and a linear head on frozen features, the labels first fix the decision thresholds: for single-modality audio and face at $K \le 5$, \NUM{p1-e29-critshare-kle5-single-lo}--\NUM{p1-e29-critshare-kle5-single-hi} of the improvement that pooling brings comes from the threshold, while the direction of audio classifiers improves only at the rate $q/K$ with $q$ between \NUM{p1-qdir-audio-lo} and \NUM{p1-qdir-audio-hi}.
A source supplies \NUM{p1-e32-n1src-audio-lo}--\NUM{p1-e32-n1src-audio-hi} labels per emotion of direction for audio, and neither a few labels nor a label-free choice of the utterances to label moves the direction of a good start much.
Weak labels on the whole unlabeled pool from a teacher of another modality, trained on the sources alone, move it: they raise every single-modality student above the five-value start, and the extended method exceeds the strongest competitor tuned under the same protocol in every configuration and with second encoders.
The two costs and the teacher's UAR advantage over the start indicate in advance where a few labels, a source, or another modality pays off.
